\section{Message schema}\label{app:messages}
\Cref{tab:messages} in \cref{sec:bus} lists the \nMsgTypes{} message types with their required fields as implemented; the bus rejects a message that lacks a required field and charges tokens at about four characters per token.

\section{Harness files}\label{app:harness}
\begin{table}[h]
\centering\scriptsize
\begin{tabular}{llp{0.5\linewidth}}
\toprule
File & Type & Retrieval keywords \\
\midrule
\texttt{agents/arm\_L.md} & agent & left arm L reach table\_left table\_center \\
\texttt{agents/arm\_R.md} & agent & right arm R reach table\_right table\_center \\
\texttt{agents/eye.md} & agent & eye camera scene report occlusion \\
\texttt{memory/episodes.md} & memory & episodic memory past failures lessons \\
\texttt{objects/book.md} & object & book textbook heavy robust \\
\texttt{objects/bottle.md} & object & bottle water liquid upright \\
\texttt{objects/laptop.md} & object & laptop heavy fragile notebook computer \\
\texttt{objects/lunchbox.md} & object & lunchbox food container \\
\texttt{rules/bag\_exclusive.md} & rule & bag opening exclusive lock one arm at a time \\
\texttt{rules/bottle\_upright.md} & rule & bottle water upright away laptop leak \\
\texttt{rules/heavy\_bottom.md} & rule & heavy bottom stack order weight mass main compartment \\
\texttt{rules/laptop\_slot.md} & rule & laptop padded slot compartment protect \\
\texttt{rules/leave\_distractors.md} & rule & leave distractor mug plate untouched not pack \\
\texttt{rules/lunchbox\_last.md} & rule & lunchbox last final order top \\
\texttt{skills/grasp.md} & skill & grasp force slip crush payload \\
\texttt{skills/handover.md} & skill & handover transfer arm center \\
\texttt{skills/place.md} & skill & place compartment orientation insert bag capacity \\
\texttt{skills/push\_aside.md} & skill & push aside clear occlusion \\
\texttt{tasks/pack\_backpack.md} & task & pack backpack school bag items instruction \\
\bottomrule
\end{tabular}

\caption{The \nHarnessFiles{} harness files with their retrieval keywords.}
\end{table}

\section{Reward and score definitions}\label{app:rewards}
\paragraph{Brain plan reward.} $R = 0.1\,r_{\text{format}} + 0.2\,r_{\text{feasible}} + 0.6\,r_{\text{rules}} + 0.1\,r_{\text{length}}$ with an optional closed-loop rollout term of weight 0.5.
$r_{\text{format}}$ checks that the output parses as a plan; $r_{\text{feasible}}$ is 1 if the static check passes and otherwise falls with the fraction of erroneous tasks; $r_{\text{rules}}$ is the valid-item fraction from a deterministic dry run minus 0.25 per rule violation and per crush, floored at 0; $r_{\text{length}}$ is 1 up to 24 tasks and falls linearly to 0 at 48.
\paragraph{Episode score.} The valid-item fraction minus 0.15 per violation, 0.3 per crushed item and 0.01 per replan.
\paragraph{Credit.} $C(i,r)=R(\text{real})-R(i,r)$; oracle gain $G=-C(i,\text{oracle})$ (\cref{sec:signals}).

\section{Upstream facts and the scorer toy case}\label{app:viki}
The scorer facts in \cref{sec:scorer} were read from the upstream VIKI-R repository by the authors: simulator-free feasibility check, \texttt{random.choice} over each asset's candidate start positions on every call, correctness as feasible and length ratio at least 0.99, reward 0.1 format plus 0.9 correctness, evaluation scripts reading the first 900 level-2 and 1000 level-1 rows, and the level-2 script reading one split file while writing to a directory with another name.
\todo{Pin the repository commit; re-run the toy case (18 of 40 unseeded draws, exhaustive average 0.5, position-agnostic plan 40 of 40) against the upstream scorer. The toy numbers are from an earlier run by the authors and could not be re-run in this build because the upstream code was not available to it, which also left the scorer tests skipped.}

\section{Environment and tests}\label{app:env}
The code base is a Python package with a symbolic world, a pose-level simulator, the loop, the harness and training scripts; the physics stack is Isaac Lab 3.0 with Newton (\texttt{mujoco-warp}), with no Isaac Sim and no PhysX, and that configuration was measured on the RTX 4050 and on a Kaggle T4 (gate G-3: 30 steps, rgb and depth contracts, D-39 and D-51), stepping one device at a time (D-50).
In the training environment (Python 3.12, torch 2.15, transformers 5.23) the test suite passes with the unavailable upstream VIKI-R clone tests skipped, so the scorer tests do not run in a clean build.
This proves that the code runs and that the wiring is consistent, and is not evidence about model quality.
\todo{Quote the current passed/skipped count from the gate record at submission time; earlier figures here, including the CPU-only 137/13 baseline, are superseded by the measured training-venv count.}

\section{Fixture mechanics checks}\label{app:fixture}
All runs here are class FIXTURE: a symbolic world, a scripted Brain, parametric Eyes and a flaky-planner wrapper, seeds 0 to \fxN{} minus one.
They check that mechanisms behave as designed and are not claims about trained agents.
The numbers come from \texttt{scripts/collect\_fixture.py}.

\begin{table}[h]
\centering\small
\begin{tabular}{lccc}
\toprule
Verification & Successes & Wilson 95\% & Retries used \\
\midrule
none & 48/80 & [0.49, 0.70] & 0 \\
static & 63/80 & [0.69, 0.86] & 29 \\
dry run & 80/80 & [0.95, 1.00] & 48 \\
dry run, 0 retries & 48/80 & [0.49, 0.70] & 0 \\
\bottomrule
\end{tabular}

\caption{Plan verification against a planner wrapped to inject errors with probability 0.5 (clean Eyes).}
\end{table}
\begin{table}[h]
\centering\small
\begin{tabular}{lcc}
\toprule
Mode & Successes & Message tokens per episode \\
\midrule
push & 73/80 & 10{,}551 \\
delta & 73/80 & 5{,}737 \\
pull & 72/80 & 1{,}895 \\
init-only & 71/80 & 1{,}182 \\
\bottomrule
\end{tabular}

\caption{Communication modes (noisy Eyes): message tokens are a four-characters-per-token proxy.}
\end{table}
\begin{table}[h]
\centering\small
\begin{tabular}{lccc}
\toprule
Condition & Successes & Wilson 95\% & Mean rounds \\
\midrule
probe & 80/80 & [0.95, 1.00] & 16.56 \\
challenge only & 39/80 & [0.38, 0.60] & 14.00 \\
no fix & 22/80 & [0.19, 0.38] & 14.35 \\
\bottomrule
\end{tabular}

\caption{Heavy Eye mass noise ($\sigma=0.45$ on both Eyes): weighing probe, challenge only, and no fix.}
\end{table}
\begin{table}[h]
\centering\small
\begin{tabular}{lcc}
\toprule
Eyes (omission fault) & LOOK off & LOOK on \\
\midrule
1 & 1/80 & 48/80 \\
2 & 15/80 & 75/80 \\
\bottomrule
\end{tabular}

\caption{Effect of the look request when each Eye omits an item with probability 0.4.}
\end{table}
\begin{table}[h]
\centering\small
\begin{tabular}{lcccccc}
\toprule
Eye fault & Mean episode score & Oracle gain G [95\% CI] & mass & class & fragility & region \\
\midrule
omission & 0.992 & 0.008 [0.002, 0.016] & 0.000 & 0.000 & 0.000 & 0.000 \\
mass\_noise & 0.826 & 0.174 [0.131, 0.223] & 0.174 & 0.000 & 0.000 & 0.000 \\
frag\_noise & 0.999 & 0.001 [0.000, 0.001] & 0.000 & 0.000 & 0.001 & 0.000 \\
misclass & 0.905 & 0.095 [0.056, 0.139] & 0.000 & 0.095 & 0.000 & 0.000 \\
region\_err & 0.954 & 0.046 [0.031, 0.062] & 0.000 & 0.000 & 0.000 & 0.046 \\
stale & 0.880 & 0.120 [0.091, 0.151] & 0.120 & 0.000 & 0.000 & 0.000 \\
\bottomrule
\end{tabular}

\caption{Oracle gain $G$ for Eye~A by injected fault type, with the breakdown by field: the estimator puts the gain in the field that carries the fault.}
\end{table}

\section{Reproducing the figures and numbers}\label{app:repro}
\begin{table}[h]
\centering\small
\begin{tabular}{lp{0.62\linewidth}}
\toprule
Item & Value \\
\midrule
Training and data seeds & $\geq 100000$ \\
Evaluation seeds & $< 100000$ \\
Stochastic draws & keyed by (seed, arm, skill, object, attempt); Eye noise by (seed, eye, time) \\
Python / simulator & 3.12 / Isaac Lab 3.0 with Newton (no Isaac Sim, no PhysX) \\
Compute & RTX 4050 (6 GB, dev) and Kaggle dual T4; G-3 PASS on both GPU classes, one device at a time (D-39, D-51) \\
Code and data release & \todo{repository, Hub and DOI links} \\
\bottomrule
\end{tabular}

\caption{Reproducibility facts.}
\end{table}
\begin{small}
\begin{verbatim}
python scripts/collect_figdata.py --project ../robot_team
python scripts/collect_fixture.py --project ../robot_team
python scripts/build_diagrams.py     # 12 figures, SVG + PDF
python scripts/audit_diagram_text.py; python scripts/audit_diagram_arrows.py
python scripts/build_paper.py        # tables, number macros, pending lists
python scripts/verify_paper.py       # numbers, citations, pending cells
pdflatex main; bibtex main; pdflatex main; pdflatex main
\end{verbatim}
\end{small}
